\ssr{ЗАКЛЮЧЕНИЕ}

Цель работы достигнута: выполнен сравнительный анализ ресурсной эффективности
алгоритмов поиска элемента в массиве~--- поиска полным перебором и трёх
реализаций бинарного поиска.

Все задачи решены:
\begin{enumerate}
	\item описана решаемая согласно варианту задача поиска элемента в массиве;
	\item разработаны алгоритм поиска полным перебором, два нерекурсивных
	алгоритма бинарного поиска (без предварительного выхода и с~предварительным
	выходом) и рекурсивный алгоритм бинарного поиска;
	\item разработанные алгоритмы описаны схемами алгоритмов;
	\item описаны требования к программному обеспечению;
	\item описаны средства реализации алгоритмов, в качестве языка
	программирования выбран Swift;
	\item разработанные алгоритмы реализованы;
	\item выполнено тестирование реализаций алгоритмов, все реализации прошли
	функциональные тесты;
	\item выполнена теоретическая оценка памяти, затрачиваемой реализацией
	каждого алгоритма;
	\item оценена временная эффективность алгоритмов в лучшем и худшем случае
	в терминах числа сравнений искомой величины с другими величинами;
	\item результаты сравнительного анализа представлены в формате таблицы;
	\item сделаны выводы из сравнительного анализа реализаций разработанных
	алгоритмов по критериям временной и ёмкостной эффективности, а также
	дополнительных требований, налагаемых алгоритмом.
\end{enumerate}

По результатам сравнительного анализа можно сделать вывод, что выбор алгоритма
поиска определяется компромиссом между числом сравнений, затратами памяти
и требованиями к входному массиву:
\begin{itemize}
	\item поиск полным перебором является наиболее универсальным, так как
	применим к произвольному массиву, однако число сравнений в худшем случае
	растёт линейно с размером массива;
	\item бинарный поиск сокращает число сравнений до логарифмического, но
	применим только к отсортированному массиву, поэтому оправдан, когда массив
	уже упорядочен или поиск выполняется многократно;
	\item предварительный выход в бинарном поиске выгоден в лучшем случае, но
	в худшем случае требует примерно вдвое больше сравнений;
	\item рекурсивная реализация бинарного поиска при той же временной
	эффективности уступает нерекурсивным по памяти из-за затрат на кадры стека.
\end{itemize}